\section{第一层：生物学奠基}

\subsection{为什么“形成”章节需要生物学}

与之相邻的意识章节已经建立了一个强实证底板：
生物系统确实呈现有组织的状态标志、稳定性结构与非平凡的转变动力学。
这一结果尚不足以直接证明日常意义上的“框架形成”，
但它改变了什么样的问题在智识上是允许提出的。
本书现在不仅可以问状态如何持续，也可以问重复耦合与条件作用如何构建出随后会改变转变成本的机制。

\begin{proposition}[对历史敏感的惯性]
当前状态的惯性不仅依赖当前刺激，还依赖与该状态相关的累积强化历史、预期奖赏景观以及以往成功重新进入该状态的路径。
\end{proposition}

\begin{discussion}
这一命题并不是从某一篇论文中直接提炼出来的定理。
它是一个以\textbf{中等}强度成立的综合性主张。
意识实证论文支持我们严肃对待有结构的状态；
因果涌现与亚稳态来源则支持把历史与系统组织看作真正的解释变量，而不只是装饰性背景。
\end{discussion}

\subsection{条件作用与整合性组织}

2025 年那篇关于基因调控网络模型中联结条件作用的 Communications Biology 论文尤其有用，
因为它给出了一个可处理的例子：重复的局部互动如何产生更强的高层组织。
对这一来源必须谨慎阅读。
它\emph{并不}表明人类学习常规在字面上就是基因调控网络。
它真正为手稿带来的，是更有限但也更有价值的三点：
\begin{enumerate}[nosep]
  \item 重复的局部耦合可以改变整体组织；
  \item 高层组织可以获得解释价值；
  \item 条件作用历史会改变后续转变的形态。
\end{enumerate}

对于框架形成而言，这意味着一个机制并不只是若干离散过去行为的简单总和。
一旦线索、奖赏、预期与环境强化开始彼此稳定，所形成的模式就会成为一个真实的运行层。
用业务语言说，系统开始在其自身结构中“携带政策”。

\subsection{亚稳态，而不是简单刚性}

2024 年 Nature Reviews Neuroscience 关于亚稳态的文章至关重要，
因为它阻止了一种错误的二元对立。
人们常把框架描述为要么僵硬的习惯，要么灵活的选择。
亚稳态提供了更好的中项：
有组织的模式可以既足够持久，又能在条件变化时重新配置。

\begin{definition}[亚稳态框架]
\textbf{亚稳态框架}是这样一种机制：它在通常扰动下具有局部持续性，
但并不会在更大的情境变化、奖赏变化或控制参数变化面前保持绝对固定。
\end{definition}

\begin{discussion}
这一定义很重要，因为许多高产框架本来就必须是亚稳态的，而不是绝对稳定的。
一种永远无法被中断的学习常规是失能的；
一种在任何扰动下都会崩塌的常规则没有价值。
因此，本章把亚稳态视为解释“好框架为什么显得稳定但又不至于变成牢笼”的正确生物学与系统层桥梁。
\end{discussion}